\section{Технологический стек}

\subsection{Серверная часть}

\headingtextsep

Серверная часть системы проектируется как модульный монолит на языке Rust. Такой выбор обусловлен требованиями к предсказуемому потреблению ресурсов, строгой типизации прикладной логики и удобству последующего выделения отдельных подсистем в самостоятельные сервисы без пересмотра базовой доменной модели.

В качестве HTTP-слоя используется фреймворк axum. Он применяется для построения прикладного API, обработки маршрутов, организации middleware и реализации точек интеграции для пользовательского интерфейса, внешних сервисов и агентного доступа.

Для асинхронного выполнения операций используется runtime tokio. Он необходим для обработки сетевых запросов, фоновых задач, операций поиска, обмена данными с внешними сервисами и работы с файловыми ресурсами.

Для журналирования и трассировки событий предусматривается использование библиотеки tracing. Она должна применяться для диагностики ошибок, отслеживания хода пользовательских и агентных операций, а также для поддержки аудита технически значимых действий.

Для работы с PostgreSQL в серверной части целесообразно использовать библиотеку sqlx, обеспечивающую типобезопасный доступ к данным и удобную интеграцию с асинхронной моделью Rust.

\textheadingsep
\subsection{Клиентская часть}

\headingtextsep

Клиентская часть системы реализуется как веб-приложение на основе React и TypeScript. Такой стек выбран для построения насыщенного интерфейса с большим количеством рабочих экранов, поддержкой сложной навигации и разделением пользовательских и административных сценариев.

React используется как основа компонентной модели интерфейса. На его базе должны быть реализованы рабочие области задач, документации, заметок, файловых ресурсов, административных разделов и экранов аутентификации.

TypeScript применяется для строгой типизации клиентской логики, контрактов взаимодействия с API и моделей представления данных. Это позволяет уменьшить количество ошибок на стыке интерфейса и серверной части.

В качестве стартового инструмента сборки выбран Vite. Он обеспечивает быстрый локальный запуск, удобную разработку пользовательского интерфейса и предсказуемую сборку веб-приложения для развертывания.

Клиентское приложение проектируется как единое веб-приложение. Пользовательская рабочая область, администрирование рабочего пространства и операторские функции на первом этапе размещаются в общей кодовой базе и разделяются маршрутами, layout-структурой и проверками прав доступа.

\textheadingsep
\subsection{Хранение данных}

\headingtextsep

Основным хранилищем структурированных данных системы является PostgreSQL. В данной СУБД предполагается хранить сущности рабочих пространств, проектов, задач, документов, заметок, участников, агентных учётных данных, связей между объектами и метаданных файловых ресурсов.

Для полнотекстового и семантического поиска используется гибридный подход на базе стандартных механизмов PostgreSQL и расширения pgvector. Полнотекстовый поиск применяется для поиска по содержимому и метаданным, а pgvector используется для поиска по смысловому сходству.

Индексирование данных должно выполняться по документам, заметкам, описаниям задач, именам и метаданным файловых ресурсов, а также по результатам агентной активности и связанным внешним артефактам. Итоговая поисковая выдача формируется как объединение результатов полнотекстового и векторного поиска.

Файловые ресурсы не хранятся непосредственно в PostgreSQL. В базе данных размещаются только метаданные файлов, связи с сущностями системы и служебная информация. Бинарное содержимое выносится в отдельный backend хранения через абстракцию AssetStorage.

На этапе локальной разработки и в простых self-hosted сценариях должен поддерживаться backend локальной файловой системы. Для production и устойчивых hosted-развертываний должен использоваться S3-совместимый backend объектного хранения.

\textheadingsep
\subsection{Инфраструктура и интеграции}

\headingtextsep

Базовой моделью развертывания системы является контейнеризированная среда. Для локального запуска и разработки предусматривается использование Docker Compose, позволяющего поднимать согласованный набор сервисов, включающий API, веб-приложение, PostgreSQL и вспомогательные сервисы.

Архитектура системы предусматривает отдельный MCP-модуль для взаимодействия с программными агентами. Этот модуль реализуется как самостоятельная Rust-утилита, работающая поверх основного API и не имеющая прямого доступа к базе данных или внутренним модулям приложения.

Для пользовательской аутентификации на первой итерации предусматривается вход через GitHub OAuth. При этом внутренняя модель идентичности проектируется независимо от конкретного провайдера, чтобы в дальнейшем можно было перейти к полноценному OIDC-провайдеру без изменения доменной модели.

Интеграция с GitHub рассматривается как основная внешняя интеграция первой версии системы. Она должна обеспечивать связывание задач с issue, pull request, commit и branch, а также обогащение проектного контекста инженерными артефактами.

Для локальной разработки предусматривается отдельный dev-only профиль аутентификации и конфигурации, допускающий запуск без внешних учётных данных. Такой профиль должен быть жёстко ограничен средой разработки и не использоваться в рабочем развёртывании.